\documentclass[11pt]{article}
\usepackage[margin=1in]{geometry}
\usepackage{graphicx}
\usepackage{booktabs}
\usepackage{xcolor}
\usepackage{hyperref}
\usepackage{amsmath}
\usepackage{listings}
\usepackage{enumitem}
\usepackage{caption}
\usepackage{float}
\usepackage{longtable}
\usepackage{xurl}
\setlength{\emergencystretch}{3em}

\hypersetup{colorlinks=true, linkcolor=blue, urlcolor=blue, citecolor=blue}

\definecolor{codebg}{rgb}{0.96,0.96,0.96}
\lstset{
  basicstyle=\ttfamily\small,
  backgroundcolor=\color{codebg},
  frame=single,
  breaklines=true,
  columns=fullflexible,
  tabsize=2,
}

% Placeholder for a screenshot that has not been captured yet.
\newcommand{\pending}[1]{%
  \begin{center}\fcolorbox{red!60!black}{red!4}{\parbox{0.92\textwidth}{\centering\vspace{1.2em}%
  \textbf{\color{red!60!black}[Screenshot pending]}\\[0.4em]#1\vspace{1.2em}}}\end{center}}

\newcommand{\shot}[3][\textwidth]{%
  \begin{figure}[H]\centering\includegraphics[width=#1]{figures/#2}\caption{#3}\label{fig:#2}\end{figure}}

\title{\Large\textbf{Lab Assignment Report: Packet Capture and\\ Traffic Analysis with Wireshark}}
\author{Sajjad Ahmed Shaaz \\ 002410501093 \\ BCSE-III A3}
\date{November 2, 2026}

\begin{document}
\maketitle

\section*{Basic Information}
\begin{itemize}[leftmargin=1.6em]
  \item \textbf{Subject:} CSE/PC/B/S/314 Computer Networks Lab
  \item \textbf{Assignment Number:} 5
  \item \textbf{Problem Statement:} Capture live network traffic with Wireshark and analyse
  it layer by layer: ARP and ICMP generated by \texttt{ping}, a complete HTTP exchange with a
  web server, the Ethernet/IP/TCP headers and raw bytes of individual packets, NIC
  manufacturers from MAC OUIs, protocol statistics, and the flow graph of the conversation.
  \item \textbf{Submission due:} first lab after the Puja vacation (tentatively 26--30
  October 2026)
\end{itemize}

\section{Design}

\subsection{Purpose}
Assignments 1--4 built protocols from scratch over sockets: error detection, flow control,
CSMA and CDMA. This assignment goes the other way. It takes the \emph{real} protocol stack
of an ordinary laptop and observes it on the wire. Every abstraction from the earlier
assignments has a concrete counterpart in the capture. Frames carry a 4-byte FCS (Assignment 1),
TCP's sliding window and retransmissions are visible (Assignment 2), and the Wi-Fi medium is
shared with other stations (Assignment 3). The goal is to identify each layer in real packets
and explain what it is doing.

\subsection{Experimental Setup}
\begin{table}[H]
\centering
\begin{tabular}{ll}
\toprule
\textbf{Item} & \textbf{Value} \\
\midrule
Host OS & Ubuntu Linux (Wireshark GUI) \\
Capture interface & \texttt{wlp0s20f3} (Intel Wi-Fi adapter) \\
Host MAC address & \texttt{f0:77:c3:ec:da:24} (\texttt{Intel\_ec:da:24}) \\
Host IPv4 / IPv6 & \texttt{10.147.196.20} / \path{2401:4900:b30f:ba3:ec7:b2ba:1982:6ec1} \\
Default gateway / DNS server & \texttt{10.147.196.35}, MAC \texttt{d2:0b:02:cf:e6:60} \\
Browser & Google Chrome 144 (User-Agent in Figure~\ref{fig:q5_follow_stream.png}) \\
HTTP test site & \path{http://httpforever.com/} (plain HTTP, no HTTPS redirect) \\
\bottomrule
\end{tabular}
\caption{Capture environment, as read from the captured packets themselves.}
\label{tab:setup}
\end{table}

\texttt{httpforever.com} was chosen because it is deliberately served over plain HTTP on
port 80 and never redirects to HTTPS. On almost every other site today the HTTP request and
response are inside TLS and cannot be read in Wireshark. Here the \texttt{GET} request,
headers and HTML body are visible in clear text. For contrast, an HTTPS site
(\path{jadavpuruniversity.in}) was also browsed. Its traffic appears only as TLS records
(\S\ref{sec:q2a}).

\subsection{Structure of the Analysis}
Every question in the assignment maps to one Wireshark view:
\begin{itemize}[leftmargin=1.6em]
  \item \textbf{Packet List} (top pane): one row per frame, used for Q1, Q2(a--c), Q5.
  \item \textbf{Packet Details} (middle pane): the protocol tree Frame $\to$ Ethernet $\to$
  IP $\to$ TCP/UDP/ICMP $\to$ application. Used for Q2(d--e), Q6, Q7.
  \item \textbf{Packet Bytes} (bottom pane): hex and ASCII of the raw frame, used for Q3, Q4, Q8.
  \item \textbf{Statistics menu}: Protocol Hierarchy (Q9) and Flow Graph (Q10).
  \item \textbf{Follow Stream}: reassembles one TCP conversation (Q5a).
\end{itemize}

\section{Procedure (Test Run)}
\begin{enumerate}[leftmargin=1.8em]
  \item Started Wireshark and opened \textbf{Capture $\to$ Options} (Ctrl+K). The Input tab
  lists all interfaces. Promiscuous mode was left enabled, and the Wi-Fi interface
  \texttt{wlp0s20f3} was selected for capture (Figure~\ref{fig:capture_options.png}).
  \item Under the \textbf{Options} tab, \emph{Update list of packets in real time} and
  \emph{Resolve MAC addresses} were enabled. MAC name resolution is what lets Wireshark
  display \texttt{Intel\_ec:da:24} instead of \texttt{f0:77:c3:ec:da:24}.
  \item Started the capture, generated traffic (ping, then browsing), and stopped it with
  Ctrl+E.
  \item Switched \textbf{View $\to$ Time Display Format $\to$ Time of Day} for the timing
  question.
\end{enumerate}

\shot[0.8\textwidth]{capture_options.png}{Capture $\to$ Options, Input tab: the list of
capture interfaces. Promiscuous mode is enabled for all interfaces.}

\shot{capture_options_tab.png}{Capture $\to$ Options, \textbf{Options} tab: \emph{Update list of
packets in real-time} and \emph{Resolve MAC addresses} are ticked. Network-name resolution is
left off, so Wireshark generates no extra DNS lookups of its own.}

\section{Questions}

%--------------------------------------------------------------------------
\subsection{Q1: ICMP Traffic with Ping, and ARP}
The gateway \texttt{10.147.196.35} was pinged from the terminal
(\texttt{ping 10.147.196.35}, 9 echo requests). Applying the display filter \texttt{icmp}
isolates the 18 ICMP packets (Figure~\ref{fig:q1_icmp_filter.png}); clearing the filter
shows the ARP exchange in between (Figure~\ref{fig:q1_ping_arp.png}).

\shot{q1_icmp_filter.png}{Display filter \texttt{icmp}: 9 Echo (ping) requests from
\texttt{10.147.196.20} and 9 replies from \texttt{10.147.196.35}. Frame 272 (first request) is
selected; its bytes start with the destination MAC \texttt{d2 0b 02 cf e6 60}, then the
source MAC \texttt{f0 77 c3 ec da 24}, then EtherType \texttt{08 00} (IPv4).}

\shot{q1_ping_arp.png}{Unfiltered view of the same period: ICMP echo pairs, with an ARP
request/reply (frames 307--308) in between.}

\paragraph{Observations.}
\begin{itemize}[leftmargin=1.6em]
  \item Each request is matched to its reply by the identifier \texttt{id=0x0001} and the
  sequence number \texttt{seq=1\ldots9}. Wireshark links them with ``reply in'' and
  ``request in''.
  \item Every frame is 98 bytes: 14 (Ethernet) + 20 (IPv4) + 8 (ICMP header) + 56 (Linux
  \texttt{ping}'s default payload). The payload bytes \texttt{10 11 12 \ldots\ 37} are
  visible in the bytes pane.
  \item The round-trip times computed from the timestamps are shown in Table~\ref{tab:rtt}.
  They range from 8.5\,ms to 23.5\,ms, with a mean of 13.3\,ms. The reply TTL is 64, the
  Linux default, so the gateway is zero routed hops away.
\end{itemize}

\begin{table}[H]
\centering
\small
\begin{tabular}{rrrrrrrrrr}
\toprule
\textbf{seq} & 1 & 2 & 3 & 4 & 5 & 6 & 7 & 8 & 9 \\
\midrule
\textbf{RTT (ms)} & 11.42 & 11.04 & 14.92 & 8.48 & 23.52 & 12.80 & 11.82 & 11.63 & 13.95 \\
\bottomrule
\end{tabular}
\caption{Ping round-trip times, reply time $-$ request time from the capture.}
\label{tab:rtt}
\end{table}

\paragraph{ARP.} Frame 307 is an ARP request \emph{``Who has 10.147.196.35? Tell
10.147.196.20''}, and frame 308 is the reply \emph{``10.147.196.35 is at
d2:0b:02:cf:e6:60''}, received 23\,ms later. Note that this request was sent to the
gateway's unicast MAC, not to the broadcast address \texttt{ff:ff:ff:ff:ff:ff}, and it
appears \emph{between} pings 6 and 7 rather than before ping 1. This is because the
gateway's MAC was already in the ARP cache when the ping started, from earlier traffic. The
first ping could therefore be sent immediately, and frame 307 is Linux's periodic
\emph{reachability probe} that re-confirms a cached entry. The classical sequence, a
\emph{broadcast} ARP request followed by the first ICMP echo, appears only when the cache is
empty.

\paragraph{Second run: ping and ARP side by side.} To look at ARP more closely, a fresh
capture was taken while running \texttt{ping -c 4 10.147.196.35}
(Figures~\ref{fig:q1_ping_new.png} and~\ref{fig:q1_arp_new.png}). All 4 echoes were answered
(0\% loss, RTT min/avg/max = 7.69/11.37/17.90\,ms; frames 103--132, identifier
\texttt{id=0x0003}). With the display filter \texttt{arp}, the capture shows three ARP
exchanges, at 15:10:13.51, 15:10:30.71 and 15:10:49.02, roughly every 17--19\,s. Each one
has the same form:
\begin{itemize}[leftmargin=1.6em]
  \item \textbf{Request} from the gateway \texttt{d2:0b:02:cf:e6:60}: \emph{``Who has
  10.147.196.20? Tell 10.147.196.35''}. It carries the sender's IP and MAC and the target
  IP.
  \item \textbf{Reply} from this laptop \texttt{Intel\_ec:da:24}: \emph{``10.147.196.20 is
  at f0:77:c3:ec:da:24''}, sent about 17\,\textmu s later.
\end{itemize}
The ping started at 15:10:15.98, only 2.5\,s after the exchange in frames 72--73. When Linux
receives an ARP request addressed to itself, it also records the \emph{sender's} IP and MAC
in its own cache. The gateway's request therefore told the laptop that 10.147.196.35 is at
\texttt{d2:0b:02:cf:e6:60}, and the laptop never had to ask. This gateway (a Wi-Fi hotspot)
keeps both caches warm with periodic unicast probes, which is why neither capture shows a
broadcast request from the PC. The ARP mechanism is identical either way: a request naming a
target IP, answered by its owner with its MAC address. Only the direction and the
broadcast/unicast destination differ.

\shot{q1_ping_new.png}{Second run: \texttt{ping -c 4 10.147.196.35} in the terminal (right)
and the four request/reply pairs in Wireshark with filter \texttt{icmp}.}
\shot{q1_arp_new.png}{Second run, filter \texttt{arp}: three request/reply exchanges, each
initiated by the gateway (\texttt{10.147.196.35}) and answered by this laptop.}

%--------------------------------------------------------------------------
\subsection{Q2: Web Traffic}

\subsubsection*{(a) Protocols in the unfiltered packet list}
\label{sec:q2a}
Figures~\ref{fig:q2_unfiltered_1.png}--\ref{fig:web_ju_2.png} show the unfiltered packet list
while browsing \texttt{httpforever.com} (HTTP) and \path{jadavpuruniversity.in} (HTTPS).
The Protocol column contains:

\begin{table}[H]
\centering
\small
\begin{tabular}{p{2.6cm}p{3.3cm}p{8.2cm}}
\toprule
\textbf{Protocol} & \textbf{Layer / transport} & \textbf{What it was doing in this capture} \\
\midrule
ARP & Link & Resolving the gateway's MAC address \\
ICMP & Network (IPv4) & Ping echo request/reply (Q1) \\
ICMPv6 & Network (IPv6) & Neighbour/router discovery (IPv6 equivalent of ARP) \\
DNS & Application / UDP 53 & Resolving \texttt{httpforever.com}, \path{jadavpuruniversity.in}, etc. \\
TCP & Transport & Handshakes, ACKs, FIN/RST, retransmissions \\
HTTP & Application / TCP 80 & \texttt{GET /}, \texttt{200 OK} from \texttt{httpforever.com} \\
HTTP/XML & Application / TCP 80 & Response carrying an XML document (\texttt{favicon.svg}) \\
TLSv1, TLSv1.2, TLSv1.3 & Application / TCP 443 & HTTPS sessions (Client Hello, Server Hello, Application Data) \\
QUIC & Transport / UDP 443 & HTTP/3 to Google servers (Chrome background traffic) \\
UDP & Transport & Datagrams not decoded further (e.g.\ QUIC short headers) \\
\bottomrule
\end{tabular}
\caption{Protocols appearing in the Protocol column. The Protocol Hierarchy statistics
(Q9) additionally list DHCP, NTP, mDNS and SSDP.}
\label{tab:protocols}
\end{table}

\shot{q2_unfiltered_1.png}{Unfiltered capture, start of the \texttt{httpforever.com} visit:
DNS queries/responses (A, AAAA and HTTPS records), QUIC, TCP SYNs to port 80 and 443, ARP.}
\shot{q2_unfiltered_2.png}{Unfiltered capture, continued: the \texttt{GET / HTTP/1.1}
(frame 106), TLS alerts, the TCP segments of the response and \texttt{HTTP/1.1 200 OK}
(frame 125), followed by requests for \texttt{style.css} and \texttt{theme.js}.}
\shot{web_ju_1.png}{Unfiltered capture while opening \texttt{https://jadavpuruniversity.in}:
DNS, the TCP three-way handshake with \texttt{182.71.247.31}, TLSv1.3 Client Hello
(SNI\,=\,\texttt{jadavpuruniversity.in}), QUIC.}
\shot{web_ju_2.png}{HTTPS continued: Server Hello and encrypted Application Data. Unlike
the HTTP site, no request or page content is visible. TCP retransmissions (frames
492--493) are highlighted in black.}

\subsubsection*{(b) Time from HTTP GET to HTTP OK}
With \textbf{View $\to$ Time Display Format $\to$ Time of Day}
(Figure~\ref{fig:q2_timeofday_get_ok.png}):

\begin{table}[H]
\centering
\begin{tabular}{rlll}
\toprule
\textbf{Frame} & \textbf{Time of day} & \textbf{Relative time (s)} & \textbf{Info} \\
\midrule
106 & 13:15:58.042404880 & 7.619590355 & \texttt{GET / HTTP/1.1} \\
125 & 13:15:58.460731909 & 8.037917384 & \texttt{HTTP/1.1 200 OK} \\
\midrule
\multicolumn{2}{r}{\textbf{Difference}} & \textbf{0.418327029\,s} & \textbf{$\approx$ 418.3\,ms} \\
\bottomrule
\end{tabular}
\end{table}

Wireshark computes the same value itself: in the details of frame 125 the HTTP layer shows
\texttt{[Time since request: 418.327029 milliseconds]} and \texttt{[Request in frame: 106]}
(Figure~\ref{fig:q2_ok_http_layer.png}). The page was gzip-compressed (3713 bytes on the
wire, 11513 bytes after decompression) and split over four TCP segments (frames 119, 120,
121 and 125). The ``200 OK'' line appears on the \emph{last} segment, because Wireshark
attaches the HTTP decode to the frame that completes the reassembled response.

\shot{q2_timeofday_get_ok.png}{Time column switched to time-of-day. Frame 106: GET at
13:15:58.042; frame 125: 200 OK at 13:15:58.460.}
\shot{q2_timeofday_dns.png}{Same capture from the beginning, in time-of-day format: DNS
lookups for \texttt{httpforever.com} at 13:15:56.61.}

\subsubsection*{(c) IP address of the website and of this computer}
\begin{itemize}[leftmargin=1.6em]
  \item \textbf{Website \texttt{httpforever.com}:} the DNS responses (frames 18--19,
  Figure~\ref{fig:q2_addresses_1.png}) returned A records \texttt{172.67.132.115} and
  \texttt{104.21.4.210}, and AAAA records \path{2606:4700:3031::6815:4d2} and
  \path{2606:4700:3036::ac43:8473}. These are Cloudflare addresses (the response header
  says \texttt{Server: cloudflare}). The browser actually fetched the page over
  \textbf{IPv6 from \path{2606:4700:3031::6815:4d2}} (frames 106/125).
  \item \textbf{This computer:} IPv4 \texttt{10.147.196.20} (private, behind NAT) and
  global IPv6 \path{2401:4900:b30f:ba3:ec7:b2ba:1982:6ec1}. The HTTP exchange used the
  IPv6 address.
\end{itemize}

\shot{q2_addresses_1.png}{Source/Destination columns widened to show the full addresses.
Frames 18--19: DNS A/AAAA answers for \texttt{httpforever.com}. Frames 26--27: TCP SYNs from
\texttt{2401:4900:b30f:ba3:ec7:b2ba:1982:6ec1} to \texttt{2606:4700:3031::6815:4d2} port 80.}
\shot{q2_addresses_2.png}{The GET (frame 106) and 200 OK (frame 125) between the same two
IPv6 addresses.}

\subsubsection*{(d) The HTTP GET packet, with the HTTP layer expanded}
Frame 106 is the \texttt{GET / HTTP/1.1} (515 bytes on the wire, TCP segment of 429 bytes,
client port 55144 $\to$ server port 80). Figure~\ref{fig:q2_get_http_layer.png} shows it
with the \emph{Hypertext Transfer Protocol} layer expanded. The request line and headers
are:
\begin{lstlisting}
GET / HTTP/1.1
Host: httpforever.com
Connection: keep-alive
Upgrade-Insecure-Requests: 1
User-Agent: Mozilla/5.0 (X11; Linux x86_64) AppleWebKit/537.36 (KHTML, like Gecko) Chrome/144.0.0.0 Safari/537.36
Accept: text/html,application/xhtml+xml,application/xml;q=0.9,image/avif,image/webp,image/apng,*/*;q=0.8,application/signed-exchange;v=b3;q=0.7
Accept-Encoding: gzip, deflate
Accept-Language: en-US,en;q=0.9,bn;q=0.8
\end{lstlisting}
The corresponding response (frame 125) with its HTTP layer expanded is shown in
Figures~\ref{fig:q2_ok_http_layer.png} and~\ref{fig:q2_ok_body.png}.

Each header ends in \texttt{\textbackslash r\textbackslash n}, and an empty line
(\texttt{\textbackslash r\textbackslash n}) ends the request. Since a GET has no body, that
blank line is the end of the message. Wireshark also links the request to its response with
\texttt{[Response in frame: 125]} and reconstructs \texttt{[Full request URI:
http://httpforever.com/]}.

\shot{q2_get_http_layer.png}{Frame 106 (\texttt{GET / HTTP/1.1}) with the HTTP layer
expanded. Note the Ethernet line: \texttt{Src: Intel\_ec:da:24}, \texttt{Dst:
d2:0b:02:cf:e6:60}, the reverse of the response in Q6.}

\shot{q2_ok_http_layer.png}{Frame 125 (response) with the HTTP layer expanded: status line,
headers (\texttt{Server: cloudflare}, \texttt{Content-Encoding: gzip}, \texttt{Transfer-Encoding:
chunked}), and Wireshark's \texttt{[Request in frame: 106]} and \texttt{[Time since request]}.}
\shot{q2_ok_body.png}{Further down the same packet: the chunked body (3713 octets),
decompressed to 11513 bytes, and the HTML (\texttt{<title>HTTP Forever}).}

\subsubsection*{(e) Value of the Host header}
\texttt{Host: httpforever.com}. HTTP/1.1 requires this header. Many sites share one IP
address (here, a Cloudflare edge server), so the server uses \texttt{Host} to decide which
site's content to return.

%--------------------------------------------------------------------------
\subsection{Q3: Hex and ASCII in the Packet Bytes Panel}
Selecting a field in Packet Details highlights the corresponding bytes in both the hex
column and the ASCII column of the Packet Bytes pane. In
Figure~\ref{fig:q3_hex_ascii.png}, the first line of the HTML body is selected in the
\emph{Uncompressed entity body} tab. Bytes \texttt{3c 21 44 4f 43 54 59 50 45 20 68 74 6d 6c
3e 0a} are highlighted in hex, alongside \texttt{<!DOCTYPE html>.} in ASCII. Each pair of
hex digits is one byte, and the ASCII column prints the character for each printable byte
(\texttt{.} otherwise). The tabs under the pane (\emph{Packet}, \emph{Reassembled TCP},
\emph{De-chunked entity body}, \emph{Uncompressed entity body}) show the same data at each
decoding stage.

\shot{q3_hex_ascii.png}{Packet Bytes panel: hex (left) and ASCII (right) with the selected
field highlighted in both.}

%--------------------------------------------------------------------------
\subsection{Q4: First 4 Bytes of the Host Parameter (Hex)}
In the GET packet the header line is the ASCII text \texttt{Host: httpforever.com\textbackslash r\textbackslash n}.
Its first four bytes, the field name \texttt{Host}, are \textbf{\texttt{48 6f 73 74}}
(H\,=\,0x48, o\,=\,0x6f, s\,=\,0x73, t\,=\,0x74). The first four bytes of the \emph{value}
\texttt{httpforever.com} are \textbf{\texttt{68 74 74 70}} (``http'').

These bytes appear in the Packet Bytes pane of frame 106 (Figure~\ref{fig:q2_get_http_layer.png}),
in the row at offset \texttt{0x0060}:
\begin{lstlisting}[columns=fixed,keepspaces=true]
0060  2f 31 2e 31 0d 0a 48 6f  73 74 3a 20 68 74 74 70   /1.1..Host: http
      /  1  .  1  \r \n H  o   s  t  :  _  h  t  t  p
\end{lstlisting}
The Host header starts at byte offset \texttt{0x66} (102). That is exactly
14 (Ethernet) + 40 (IPv6) + 32 (TCP with timestamp option) = 86 bytes of headers, plus the
16 bytes of \texttt{GET / HTTP/1.1\textbackslash r\textbackslash n}.

\shot{q4_host_bytes.png}{Enlarged Packet Bytes pane of frame 106 (from
Figure~\ref{fig:q2_get_http_layer.png}). The red boxes were added to mark the bytes
\texttt{48 6f 73 74} and their ASCII form \texttt{Host}.}

%--------------------------------------------------------------------------
\subsection{Q5: Filtering by Protocol and Following a Stream}
Display filters were applied for each protocol:
\begin{itemize}[leftmargin=1.6em]
  \item \texttt{icmp}: Figure~\ref{fig:q1_icmp_filter.png} (18 packets).
  \item \texttt{http}: Figure~\ref{fig:q5_http_filter.png}. 17 HTTP packets: the
  \texttt{httpforever.com} page and its sub-resources (\texttt{style.css},
  \texttt{theme.js}, \texttt{favicon.svg}), plus periodic \texttt{GET /} requests to
  \texttt{91.189.91.x} / \texttt{2620:2d:4002:1::x} answered with \texttt{204 No
  Content}. Those are Ubuntu NetworkManager's connectivity check
  (\path{connectivity-check.ubuntu.com}, see DNS frames 1--4), which tests for a captive
  portal.
  \item \texttt{tcp}: Figure~\ref{fig:q5_tcp_filter.png}, 521 of 1519 packets (34.3\%) in the
  saved capture. It shows the handshakes (\texttt{SYN}, \texttt{SYN,ACK}, \texttt{ACK}), the
  TLS Client Hello attempts, and the \texttt{RST}s Chrome sends to abandon the losing
  IPv4 connections (red rows; see Q10).
  \item \texttt{dns}: Figure~\ref{fig:q5_dns_filter.png}, 146 packets (9.6\%). Every query
  goes from \texttt{10.147.196.20} to the resolver \texttt{10.147.196.35} on UDP port 53.
  For each name the browser asks three questions at once: \texttt{A} (IPv4),
  \texttt{AAAA} (IPv6) and \texttt{HTTPS} (type 65, which advertises HTTP/3 and HTTPS
  support). Frames 12--20 are the lookups for \texttt{httpforever.com}.
\end{itemize}

\shot{q5_http_filter.png}{Display filter \texttt{http}.}

\shot{q5_tcp_filter.png}{Display filter \texttt{tcp}.}
\shot{q5_dns_filter.png}{Display filter \texttt{dns}.}

\paragraph{(a) Following a conversation.} Right-clicking frame 125 and choosing
\textbf{Follow $\to$ TCP Stream} applies the filter \texttt{tcp.stream eq 1} and shows the
whole conversation reassembled. The client's request is in red and the server's response in
blue (Figure~\ref{fig:q5_follow_stream.png}). The stream contains the GET for \texttt{/},
the full response headers, and the decompressed HTML of the page. Wireshark reports
``2 client pkts, 2 server pkts, 3 turns'' and 25\,kB in total. Because this is plain HTTP,
everything the browser sent and received is readable. Following a TLS stream from the
\path{jadavpuruniversity.in} visit would show only encrypted bytes.

\shot{q5_follow_stream.png}{Follow TCP Stream on stream 1 (\texttt{httpforever.com}, client
port 55144 $\leftrightarrow$ server port 80).}

%--------------------------------------------------------------------------
\subsection{Q6: HTTP Packet from the Server (TCP Source Port 80), Ethernet Layer}
Frame 125 (\texttt{HTTP/1.1 200 OK}) has \texttt{Src Port: 80}, \texttt{Dst Port: 55144}.
Its expanded Ethernet II layer (Figures~\ref{fig:q6_ethernet_src.png}
and~\ref{fig:q6_ethernet_ip.png}):
\begin{itemize}[leftmargin=1.6em]
  \item \textbf{Destination:} \texttt{Intel\_ec:da:24 (f0:77:c3:ec:da:24)}, this laptop.
  LG bit = 0: \emph{globally unique address (factory default)}. IG bit = 0: unicast.
  \item \textbf{Source:} \texttt{d2:0b:02:cf:e6:60}. LG bit = 1: \emph{locally
  administered address (this is NOT the factory default)}. IG bit = 0: unicast.
  \item \textbf{Type:} \texttt{0x86dd} (IPv6).
\end{itemize}
The IPv6 header shows \texttt{Hop Limit: 53}. If the server started at the common default of
64, the packet crossed about 11 routers. The source MAC therefore cannot be the web server's own NIC. Each
router rewrites the Ethernet header, and the source MAC of a frame always belongs to the
\emph{last hop}, which here is the gateway \texttt{10.147.196.35}. This is the same MAC that
answered the ARP in Q1. The request in the opposite direction (frame 106,
Figure~\ref{fig:q2_get_http_layer.png}) has the two addresses swapped.

\shot{q6_ethernet_src.png}{Frame 125, Ethernet layer expanded with the Source field
selected. The source MAC \texttt{d2 0b 02 cf e6 60} is highlighted at bytes 6--11.}
\shot{q6_ethernet_ip.png}{Frame 125: Ethernet and IPv6 layers expanded. Hop Limit 53,
source \texttt{2606:4700:3031::6815:4d2}.}

%--------------------------------------------------------------------------
\subsection{Q7: Manufacturers of the NICs}
\begin{itemize}[leftmargin=1.6em]
  \item \textbf{This PC's NIC: Intel.} With MAC name resolution enabled, Wireshark
  replaces the OUI \texttt{f0:77:c3} with \texttt{Intel} (\texttt{Intel\_ec:da:24}). This
  matches the laptop's Intel Wi-Fi adapter (\texttt{wlp0s20f3}).
  \item \textbf{The ``server'' side NIC: no manufacturer.} As explained in Q6, the frame's
  source MAC is that of the first-hop gateway, not of the Cloudflare server. Ethernet
  addresses are only meaningful within one link, so the server's real NIC is never visible
  from here. The gateway's MAC \texttt{d2:0b:02:cf:e6:60} has the locally-administered bit
  set, and Wireshark prints it raw with no vendor name. Such addresses are not assigned
  through the IEEE OUI registry. They are typically generated by the device's software for
  privacy (MAC randomisation, as used by Android/iOS phones acting as Wi-Fi hotspots), so no
  manufacturer can be inferred.
\end{itemize}

%--------------------------------------------------------------------------
\subsection{Q8: Hex Values of the Two OUIs}
The OUI is the first 3 bytes of a MAC address. In frame 125's raw bytes
(Figure~\ref{fig:q6_ethernet_src.png}), the first row reads:
\begin{lstlisting}[columns=fixed,keepspaces=true]
0000  f0 77 c3 ec da 24 d2 0b 02 cf e6 60 86 dd 60 0d
      \______/          \______/          \___/
      dst OUI           src OUI           EtherType (IPv6)
\end{lstlisting}
\begin{itemize}[leftmargin=1.6em]
  \item PC's NIC (destination, bytes 0--2): \textbf{\texttt{F0 77 C3}}, Intel Corporate.
  \item Gateway/server-side NIC (source, bytes 6--8): \textbf{\texttt{D2 0B 02}}. Binary of
  the first byte: \texttt{1101 00\underline{1}0}. Bit 1 (the U/L bit) is 1, so the address
  is locally administered and has no registered manufacturer.
\end{itemize}

%--------------------------------------------------------------------------
\subsection{Q9: Protocol Statistics}
\textbf{Statistics $\to$ Protocol Hierarchy} (no display filter, 1496 frames,
Figure~\ref{fig:q9_protocol_hierarchy.png}):

\begin{table}[H]
\centering
\begin{tabular}{lrrr}
\toprule
\textbf{Protocol} & \textbf{Over IPv6} & \textbf{Over IPv4} & \textbf{Total} \\
\midrule
TCP & 453 (30.3\%) & 64 (4.3\%) & \textbf{517 packets = 34.6\%} \\
UDP & 630 (42.1\%) & 185 (12.4\%) & \textbf{815 packets = 54.5\%} \\
ICMPv6 & 102 (6.8\%) & -- & 102 = 6.8\% \\
ARP & \multicolumn{2}{c}{(link layer)} & 62 = 4.1\% \\
\bottomrule
\end{tabular}
\caption{Share of packets by transport protocol (percent of all 1496 frames).}
\end{table}

\begin{itemize}[leftmargin=1.6em]
  \item \textbf{(a) TCP: 34.6\%} of packets. Higher-level protocols carried over TCP in this
  capture: \textbf{HTTP} (port 80, the \texttt{httpforever.com} page) and \textbf{TLS}
  (port 443, HTTPS).
  \item \textbf{(b) UDP: 54.5\%} of packets. Higher-level protocols carried over UDP:
  \textbf{DNS} (port 53), \textbf{QUIC} (UDP 443, HTTP/3, the single largest protocol at
  39.8\% of packets and 61.2\% of bytes), DHCP, NTP, mDNS and SSDP.
\end{itemize}
UDP outnumbering TCP may be unexpected for ``web browsing''. Chrome talks to Google services
over QUIC, which runs on UDP, and DNS, NTP, mDNS and SSDP background chatter is all UDP as
well. The roughly 79\% / 17\% split between IPv6 and IPv4 also shows that this network has
working IPv6, which the browser prefers.

\shot{q9_protocol_hierarchy.png}{Statistics $\to$ Protocol Hierarchy.}

%--------------------------------------------------------------------------
\subsection{Q10: Flow Graph}
\textbf{Statistics $\to$ Flow Graph} draws each host as a vertical line and each packet as
an arrow between them, labelled with ports and a summary. In current Wireshark the old
\emph{General Flow / Network Source} options correspond to \emph{Flow type} and
\emph{Addresses}.

\shot{q10_flow_http.png}{Flow graph limited to the \texttt{http} display filter. Top: the
\texttt{httpforever.com} page load (GET / $\to$ 200 OK, then style.css, theme.js,
favicon.svg over two parallel connections, client ports 55144 and 55152). Below: the
periodic Ubuntu connectivity checks (\texttt{GET /} $\to$ \texttt{204 No Content}), the
last one involving a TCP retransmission.}

\shot{q10_flow_all.png}{Flow graph of all TCP flows around the page load. It shows Chrome
racing several connections to \texttt{httpforever.com} at the same time, over IPv6 and IPv4
and on ports 80 and 443.}

\paragraph{What the full flow graph reveals.} Within 0.6\,s of the DNS answer, Chrome opened
TCP connections to \texttt{httpforever.com} over \emph{both} IPv6 and IPv4
(\texttt{104.21.4.210}) and on \emph{both} port 80 and port 443. That is two port-80
connections per address family, plus several port-443 attempts.
\begin{itemize}[leftmargin=1.6em]
  \item \textbf{Happy Eyeballs (RFC 8305).} IPv6 and IPv4 connections are attempted in
  parallel. The IPv6 SYN/ACKs arrived first (13:15:57.246), so the IPv4 connections were
  closed with \texttt{RST} as soon as their SYN/ACKs arrived (13:15:57.531).
  \item \textbf{HTTPS-first probing.} The port-443 connections send a TLS Client Hello
  for \texttt{httpforever.com} (in the SNI field). The server answers with a TLS \emph{Alert (Fatal,
  Handshake Failure)} because the site deliberately has no HTTPS, and Chrome falls back to
  plain HTTP on port 80.
  \item The page is then fetched over the IPv6 port-80 connection, with the second
  connection used for parallel sub-resource downloads.
\end{itemize}

\section{Analysis}

\subsection{Layering seen in a single packet}
Frame 125 alone shows the whole stack (``Protocols in frame:
\path{eth:ethertype:ipv6:tcp:http:data-text-lines}'', Figure~\ref{fig:q2_ok_frame.png}).
Ethernet delivers the frame across one hop from the gateway to the laptop, using MAC
addresses that change at every router. IPv6 carries it end to end, using addresses that stay
the same from Cloudflare to the laptop, with a hop limit that drops by one per router. TCP
ports demultiplex it to the right socket (80 $\to$ 55144), and its sequence numbers place
these 5 bytes after the previous 5204. HTTP finally interprets the reassembled 5208 bytes as
a chunked, gzip-compressed HTML response. Each layer only reads its own header, which is
exactly the encapsulation model.

\shot{q2_ok_frame.png}{Frame layer of frame 125: arrival time, frame length, and the
protocol stack \texttt{eth:ethertype:ipv6:tcp:http:data-text-lines}.}

\subsection{HTTP vs.\ HTTPS on the wire}
For \texttt{httpforever.com}, Follow TCP Stream exposed the complete request (including the
browser and OS in \texttt{User-Agent}) and the full HTML page. Anyone on the same Wi-Fi with
Wireshark could read or modify it. For \path{jadavpuruniversity.in}, the only cleartext
items are the DNS query and the server name in the TLS Client Hello's SNI field. Everything
after the Server Hello is opaque \texttt{Application Data}. This is the practical meaning
of TLS: it hides \emph{content}, but not \emph{who} you are talking to.

\subsection{Reliability mechanisms from Assignment 2, in production}
The capture contains real instances of TCP's error and flow control. Examples include
\texttt{[TCP Retransmission]} (frames 492--493 of the HTTPS capture), \texttt{[TCP Dup ACK]},
\texttt{[TCP Previous segment not captured]}, and \texttt{[TCP ACKed unseen segment]}.
Window-size fields (\texttt{Win=64896}, \texttt{Win=131072}) and SACK permission in every SYN
correspond to the Go-Back-N / Selective Repeat ideas implemented in Assignment 2. Production
TCP uses cumulative ACKs with selective acknowledgements, and its window is advertised
dynamically rather than fixed.

\subsection{Background traffic}
Much of the capture was not caused by the browsing at all. It includes NetworkManager's
captive-portal checks (HTTP 204), NTP time sync, mDNS/SSDP service discovery on the LAN,
DHCP, IPv6 neighbour discovery, and Chrome's QUIC connections to Google. Display filters
are therefore essential: the 2 interesting HTTP packets (GET and 200 OK) are among 1496 in
total.

\section{Comments}

\paragraph{Evaluation of the Lab Assignment.} This was the most directly practical of the
assignments so far. Instead of building a simplified protocol, it showed how the real ones
behave. Several things only become clear when seen on the wire: MAC addresses are
rewritten at every hop while IP addresses persist, a browser silently races several connections
to load one page, and a phone hotspot uses a randomised MAC that defeats OUI lookup.

\paragraph{Difficulty.} The tooling itself is easy. The questions, however, assume a
simpler network than most students have today. On a modern network the browser prefers
IPv6, uses QUIC for many sites, and almost every site is HTTPS-only. Some answers also differ
from what the questions expect. The ARP request was a unicast probe rather than a broadcast,
because the hotspot's own periodic ARP probes kept the laptop's cache populated. The ``server's NIC manufacturer'' cannot be
determined, because the server is about 11 hops away. Working out \emph{why} was the most
instructive part of the assignment.

\paragraph{Suggested Improvements.} Specify a plain-HTTP test site (such as
\texttt{httpforever.com}) in the assignment text. Ask students to flush the ARP cache and ping
in one command (\texttt{sudo ip neigh flush dev <iface> \&\& ping \ldots}), so the
broadcast request is captured before anything refills the cache. Reword Q7 to ask for the \emph{next-hop
router's} NIC, which is what the Ethernet header actually shows.

\end{document}
